\section{Introduction}
\label{sec:introduction}

Optimization research spans mathematical programming, gradient and derivative-free methods, evolutionary and swarm search, distribution-based optimization, surrogate methods, multiobjective optimization, robust and stochastic formulations, dynamic and online optimization, and learning-assisted approaches. At the same time, the number of named metaheuristics has grown rapidly. Large reviews have catalogued hundreds of methods, while recent structural studies have shown that differently named algorithms can share substantial computational structure. Functional taxonomies and benchmark-driven reviews have further shifted attention from inspiration metaphors toward search dynamics and empirical behavior.

These developments expose a second problem beyond algorithm proliferation: evidence is fragmented across incompatible levels. A canonical paper may define an update rule, a theoretical study may establish convergence under specific assumptions, a benchmark paper may report finite-budget behavior, a later variant may introduce adaptation, and an implementation may differ from the published specification. Broad statements about scalability or robustness can therefore mix theory, implementation, benchmark distribution, parameter policy, and computational budget unless their conditions are retained.

This paper addresses that evidence problem through a mechanism--evidence--gap atlas. Algorithm identity, computational mechanism, empirical evidence, reproducibility, and unresolved questions are represented separately before synthesis. An algorithm name is retained for provenance; a mechanism fingerprint describes search state, information, proposal geometry, memory, selection, adaptation, diversity, constraints, and models; an evidence state records what is known under a declared condition; and an evidence obligation states what observation would resolve an uncertainty.

The framework is cross-paradigm. The census contains 90 optimization identities spanning classical, derivative-free, evolutionary, swarm, probabilistic, surrogate, multiobjective, constrained, robust, dynamic, mixed-variable, and learning-assisted optimization. These identities are not interpreted as 90 distinct mechanisms. A representative evidence layer contains 42 source-backed mechanism fingerprints linked to 46 verified source records. This separation preserves broad literature coverage without equating naming diversity with mechanistic novelty.

The paper also distinguishes evidence absence from algorithmic failure. Condition-level states are represented as documented, partial, contradictory, untested, or project-validated. Structural similarity can generate a hypothesis but cannot transfer empirical evidence between algorithms. A missing experiment creates an evidence obligation rather than proof of a defect.

To connect synthesis with executable evidence, the atlas includes a frozen continuous black-box characterization using the noiseless COCO BBOB suite. Random Search, SciPy differential evolution, bounded SciPy Nelder--Mead, and pycma CMA-ES are evaluated under normalized objective-evaluation budgets across dimensions 5, 10, 20, and 40. Discovery and held-out instance sets are separated, objective-call accounting is cross-checked against COCO, and official \texttt{cocopp} postprocessing is retained for target-runtime analysis. The experiment demonstrates condition-level evidence generation rather than a universal optimizer ranking.

An important outcome is methodological restraint. The held-out campaign executed successfully, but the exact numerical reliability threshold and target subset required for a binary failure-frontier claim were not frozen before confirmation. The paper therefore retains the campaign as project-validated characterization and does not construct an exact frontier afterward. No new optimizer is introduced because the frozen evidence does not establish the independently supported mechanism-level repair opportunity required to justify one.

The contributions are fourfold. First, the paper develops a mechanism-first representation separating algorithm identity from state, information, proposal, memory, selection, adaptation, and problem-structure mechanisms. Second, it organizes heterogeneous literature into condition-level evidence states linked to provenance. Third, it couples the review to an executable reproducibility chain from protocol and configuration through frozen artifacts and claims. Fourth, it converts unresolved observations into explicit evidence obligations and a researcher decision framework for comparator selection, benchmark design, novelty checking, and future experiments.

The objective is not to replace specialized optimization theory or domain-specific benchmarking. It is to provide a common evidence architecture through which those results can be compared without erasing their assumptions and conditions. The remainder of the paper defines the review methodology, develops the mechanism taxonomy, constructs the evidence atlas, separates theory from finite-budget evidence, presents the controlled computational characterization, evaluates gaps and existing remedies, and concludes with a decision framework and bounded research frontier.
